\scanpage{8}
Ceci posé, je dis que l'hom précédent
\[
Rp_*(W)\longrightarrow e_*(\underline{\mathcal O}_S)\otimes\tau_{B^*}[-n]
\]
est un isomorphisme. Il suffit pour ceci de prouver que l'homomorphisme induit \deleted{\unsure{sur les|pour les}} en tout point $\unsure{x|b|\deleted{x}b}\in B^*$ par $\mathbb{L} i_x^*$ est un isomorphisme (histoire de mapping-cylinders\ldots). En les points $b\in B^*-e\unsure{_1(S)|(S)|_1(s)}$, ce n'est autre que le lemme 2. En un point $\deleted{b\in e}$ $b$ de la forme $e(s)$, on est ramené au cas où $S$ est le spectre d'un corps algébriquement clos, à faire une vérification de compatibilité\ldots.

\uline{Remarques} \quad 1) a) $\Phi_L$ transforme $e_A$ en $1_B$, $e_s$ ($s$ section de $A/S$) en $L_s$ (faisceau inv. alg. $\sim0$ sur $B$ rel $/S$), \inserted{donc} $e_s*F$ en $L_s\otimes\Phi_L(F)$.

b) $\Phi_L$ transforme \deleted{$\omega_A[n]$ en $\unsure{\varepsilon|e}$, $\omega_A[n]$} $L_t$ ($t$ section de $B$ sur $S$) en $\tau[-n]e_{-t}$ i.e $L_t\,\omega[n]$ en $e_{-t}$, \deleted{\unsure{de|donc}} $F\otimes L_t\deleted{\unsure{\omega[n]}}$ en $\left\{\Phi_L(F)*e_{-t}\right\}\deleted{\unsure{\omega[n]}}\unsure{e|2|,|.}$

c) $\varepsilon_{\unsure{B}}(\Phi_L(F))=\chi_A(F)$, $\chi_B(\Phi_L(F))=(-1)^n\varepsilon_A(F)$
% slanted
\margin{$\varepsilon={}$ augmentation\quad égalité \deleted{\unsure{des|dans} \unsure{fonctions}} $H^0(S,\mathbb Z)$}

2) $\det^*\circ\Phi_L$ définit \deleted{, pour tte section}
\[
\varphi\colon\underline{\mathrm{Pic}}_{A/S}\longrightarrow\underline{\mathrm{Pic}}_{B/S}
\]
qui, pour toute section $\delta$ de $\underline{\mathrm{NS}}_{A/S}$, définit
\[
\underline{\mathrm{Pic}}^\delta_{A/S}\longrightarrow\underline{\mathrm{Pic}}^{-\delta'}_{B/S}\qquad\text{compatible avec}
\]
\[
-(\widetilde{-\delta'})=\widetilde{\delta'}
\]
% archive sheet number: 7
